\section*{摘要}\label{ux6458ux8981}

早期眼底筛查是眼科疾病诊断的重要依据。围绕眼底筛查，传统机器学习、视觉大语言模型和卷积神经网络各具特点。考虑到特征建模能力、数据需求量和部署成本，本文采用卷积神经网络作为基础路线。然而，真实临床应用仍面临多疾病共存导致误检与漏检、模型跨机构泛化性能下降，以及隐私泄露风险等挑战。受以上启发，我们提出自适应标签-特征关联性构建模块，通过Transformer机制捕捉疾病共现依赖，减少误检与漏检；设计双分支引导域不变特征提取模块，在多域数据协同训练受限的条件下，仅用单一源域数据并结合CAM引导训练，提取域不变特征，提升对未知域的泛化能力，同时利用CAM可视化模型关注区域，增强可解释性；集成自适应高斯差分隐私于模型训练过程，通过动态梯度噪声注入，在中等隐私保护下兼顾模型性能。实验部分，通过完整消融验证了各模块的有效性，并使用三个数据集在域内与跨域泛化中取得性能提升，在隐私保护下性能衰减可控。综合以上，本文搭建了一个面向多标签眼科疾病识别的可解释与隐私保护单域泛化框架，提升了多疾病识别任务的跨域泛化性能，并为隐私保护场景下的眼科疾病诊断提供启发。

\noindent\textbf{关键词：}\textbf{Domain Generalization; Multi-Label Ocular Disease Recognition; Data Privacy; Interpretability}
